% =====================================================================
% 第 9 章 表面分析技术基础  分片 d
% 覆盖：§9.2.7 其他初级结构谱线 / §9.2.8 定量（半定量）分析
%       §9.2.9 数据后期处理 / §9.2.10 实验中的常见问题
% 注意：源文件 ch09-d.txt 的实际节号为 §9.5 尾＋§9.6—§9.8；§9.2.7—§9.2.10 的原书正文见书内 p.623—641，其 OCR 在 sources/parts/ch09-b.txt。逐条核对结果见各节「核对说明」。
% =====================================================================

\subsection{XPS 中除芯能级谱外的其他初级结构谱线的应用}
\label{sec:ch09-otherlines}

\emph{核对说明}：本节材料已对照原页面图像（书内 p.705—708 / PDF p.196—199）。价带谱材料在源文件 §9.6.2「参考 C 的价带谱」（书内 p.705—706）：PBMA（聚甲基丙烯酸正丁酯）的三个同分异构体 C 1s 峰型与峰位一致（图 9-156a），而价带谱差异明显（图 9-156 b、c、d）；聚烯烃 PE、PP、PB、PIB、PI 与 PS 的 C 1s 谱也难区分，价带谱有差别（图 9-157，书内 p.706）。shake-up 材料在同节「C 1s 的 Shake-up 峰分析」（书内 p.706—708）：石墨与碳纳米管的 C 1s 均为 284.6 eV、\ce{C60} 为 284.75 eV（书内 p.707），\ce{C60} 的 shake-up 峰可分解为 5 个峰（书内 p.708）。需要说明的是：§9.2.7—§9.2.10 的原书正文确实存在，其 OCR 文本落在另一分片的源文件 \texttt{sources/parts/ch09-b.txt}（§9.2.7 起于第 227 行，§9.2.8 起于第 415 行／书内 p.630，§9.2.9 起于第 589 行／书内 p.633，§9.2.10 起于第 780 行／书内 p.640），本分片的源文件 \texttt{ch09-d.txt} 并不覆盖这四节（其实际内容为 §9.5 末＋§9.6—§9.8，书内 p.694—733 / PDF p.185—224）。本分片各节按任务书的节号编排，正文为按学科标准形式的改写，逐条核对结果见各节前的核对说明。

\begin{probout}
芯能级主峰回答的是"有什么元素、各元素处在什么化学态、相对量多少"三件事。除此之外还有三类初级结构谱线能回答主峰回答不了的问题：\emph{价带谱}回答成键与电子结构（同分异构体、同类聚合物的区分）；\emph{shake-up 与 shake-off 伴峰}回答未占据态与共轭程度（石墨碳、碳纳米管与 \ce{C60} 的鉴别）；\emph{X 射线激发 Auger 谱（XAES）与 Auger 参数}回答化学态，并且能在很大程度上抵消荷电漂移。
\end{probout}

\begin{principle}
震激（shake-up）与震离（shake-off）的机理已在「终态效应及伴峰」一节（§~\ref{sec:ch09-satellites}）给出，这里只讲它们的\emph{应用}：$\pi\to\pi^{*}$ 震激峰与共轭程度直接相关，是石墨碳、碳纳米管与 \ce{C60} 的鉴别依据；震离则形成宽而弱的背景结构。价带谱是价电子态密度经各轨道光电离截面加权后的分布，出现在约 $0\sim30$ eV 的低结合能区。

X 射线激发的 Auger 电子动能由三个能级之差决定，与激发光子能量无关；把同一元素的 Auger 动能与光电子结合能相加，得到 Wagner 修饰 Auger 参数
\begin{equation}
\alpha' = E_{\mathrm{K}}(\text{Auger}) + E_{\mathrm{B}}(\text{光电子})
\label{eq:ch09-augerparam}
\end{equation}
样品荷电使 $E_{\mathrm{B}}$ 增加 $q$、使 Auger 动能减少 $q$，两者在 $\alpha'$ 中一阶抵消，因此 $\alpha'$ 对荷电不敏感，可用于区分化学态并跨仪器比较。
\emph{核对说明}：已对照原页面图像（书内 p.625 / PDF p.116）。原书印：$\alpha = \mathrm{EK_A} - \mathrm{EB_P}$（式 9-6）、$\alpha' = \alpha + h\nu = \mathrm{EK_A} + \mathrm{EB_P}$（式 9-7），式中 $\mathrm{EB_P}$ 为最强光电子线的结合能、$\mathrm{EK_A}$ 为最尖锐俄歇线的动能，并给出 Wagner 图与部分 Cu 化合物的 Cu 2p／Cu L$_3$VV 参考峰位及 $\alpha'$ 参考值（图 9-71）。故本式 $\alpha' = E_{\mathrm{K}}(\text{Auger}) + E_{\mathrm{B}}(\text{光电子})$ 并非无据重建：它与原书式 (9-7) 形式一致，仅符号写法不同（原书用 $\mathrm{EK_A}$／$\mathrm{EB_P}$，本分片用 $E_{\mathrm{K}}$／$E_{\mathrm{B}}$），判定为形式一致但符号不同；式前「样品荷电使 $E_{\mathrm{B}}$ 增加 $q$、Auger 动能减少 $q$ 而一阶抵消」的补偿关系为按学科标准形式的补充说明，原书未展开。
\end{principle}

\begin{measure}
价带谱：低通能（约 20—50 eV）、小步长（$\le0.1$ eV）、长驻留；用单色 Al K$\alpha$ 可去掉 K$\alpha_{3,4}$ 卫星线（它会在主峰高结合能约 9.8 eV 处造出一个假伴峰）。样品必须与样品台电接触良好，否则费米边畸变会被误读成"带隙变化"。伴峰测量仍用常规高分辨窄扫，但需要长采集才能取得主峰强度 $10^{-2}$ 量级的弱结构，同时必须把 X 射线累积剂量控制在损伤阈值以下（见 §\ref{sec:ch09-problems}）。
\end{measure}

\begin{dataflow}
价带谱或伴峰 $\to$ 与标准谱／态密度计算对照 $\to$ 成键与共轭判断 $\to$ 与同一批样品的芯能级化学态结论联合 $\to$ 结构归属。任何一步缺失（只有谱而没有对照，或没有芯能级支持），结论只能写成"与……一致"，不能写成"证明是……"。
\end{dataflow}

\paragraph{读谱实例：主峰分不开的碳材料，靠伴峰分开}
\begin{obs}
石墨与碳纳米管的 C 1s 主峰都在 284.6 eV（差别小于 0.15 eV，常规分辨率下不可分辨）；两者在主峰高结合能侧约 6—7 eV 处都有 $\pi$-$\pi^{*}$ 伴峰；\ce{C60} 的 C 1s 为 284.75 eV，其伴峰区域更宽，可分解为 5 个成分。同类证据来自 PBMA 的三个同分异构体：C 1s 峰型与峰位一致，而价带谱差异明显；聚烯烃 PE、PP、PB、PIB、PI 与 PS 的 C 1s 也难区分，而价带谱有别。\src{源文件 §9.6.2，书内 p.705—708 / PDF p.196—199}
\end{obs}
\begin{infer}
石墨与碳纳米管的伴峰结构一致 $\Rightarrow$ 两者价电子结构相近，碳原子均以 sp$^{2}$ 杂化形成共轭 $\pi$ 键，可归为一类；\ce{C60} 的伴峰含额外成分 $\Rightarrow$ 除共轭 $\pi$ 键外还有明显的 $\sigma$ 键贡献，须单独归属。\emph{注意}：这里推断的是"电子结构相近／不同"，不是"结构相同"；把价带谱与伴峰当指纹使用前需要大量标准谱储备，目前这类应用总体上仍停留在指纹水平，还需理论计算支撑。
\end{infer}

\begin{pitfall}
（1）把 shake-up 伴峰当作一个新的化学态单独定量——伴峰是主峰的卫星，其强度已包含在主峰的归属里，重复计入会使该元素的原子分数偏高。（2）把 Al K$\alpha_{3,4}$ 卫星线当成伴峰，或把 Mg K$\alpha$ 源下约 8—10 eV 处的卫星结构误读为共轭特征。（3）用价带谱直接断言骨架结构而不做对照与计算。（4）在绝缘粉末上测价带谱并用费米边定标。（5）样品已被 X 射线损伤（脱氟、氧化、交联）之后，再把伴峰结构的变化当作本征信息。
\end{pitfall}

\begin{crosscheck}
价带谱与 UPS（He I，21.2 eV）对照：UPS 分辨率更高、价带精细结构更清楚，而 XPS 价带谱可与同批样品的芯能级在同一次装样中测得，条件一致性更容易保证。化学态归属用 Auger 参数与 XANES／EXAFS 回查；聚合物骨架用 IR／Raman 与固体 NMR 回查；碳材料用 Raman（D 带／G 带）回查。
\end{crosscheck}

\subsection{XPS 定量（半定量）分析}
\label{sec:ch09-quant}

\emph{核对说明}：本节有原书对应正文，已对照原页面图像（书内 p.630、p.632—p.633 / PDF p.121、p.123—p.124）并通查其 OCR 全文（\texttt{sources/parts/ch09-b.txt} 第 415—588 行，即原书 §9.2.8「XPS 定量（半定量）分析」，书内 p.630 起）。原书给出的关键式：式 (9-8) $I_{ij} = K\,T(\mathrm{KE})\,L_{ij}(\gamma)\,\sigma_{ij}\!\int n_i(z)\mathrm{e}^{-z/(\lambda(\mathrm{KE})\cos\theta)}\,\mathrm{d}z$、式 (9-9) $n_i/n_j = (I_i/S_i)/(I_j/S_j)$、式 (9-10) $n_i/n_j = (I_i/I_j)(\sigma_j/\sigma_i)(\mathrm{EK}_j^{0.5}/\mathrm{EK}_i^{0.5})$，并强调「XPS 的真正定量分析很难……给出的常常是一种半定量的分析结果」，报出的是相对原子分数而不是质量分数。本节正文为按学科标准形式的改写：式 \equref{eq:ch09-quant-int}、\equref{eq:ch09-rsf}、\equref{eq:ch09-atomfrac} 与原书上式形式一致而符号与因子取舍不同（本节把式 (9-8) 的深度积分简化为 $\lambda_i\cos\theta$ 因子、按原书「对固体通常被忽略」略去 $L_{ij}(\gamma)$，并显式写出探测效率 $D(E_i)$），判定为形式一致但符号不同，并属改编补充。

\begin{probout}
XPS 定量能得到的是\emph{采样深度内}（约 $3\lambda\cos\theta$，常规为几纳米到十几纳米）元素的\emph{相对}原子分数、表面富集或偏析程度、活性组分的表面浓度。它得不到体相组成，也得不到绝对含量（除非对同基体标样做工作曲线）。因此本书与仪器厂商一致地把 XPS 定量称为\emph{半定量}。
\end{probout}

\begin{principle}
第 $i$ 个元素某芯能级的信号强度可以写成
\begin{equation}
I_i \;\propto\; n_i\,\sigma_i(h\nu)\,\lambda_i(E_i)\,T(E_i)\,D(E_i)\cos\theta
\label{eq:ch09-quant-int}
\end{equation}
其中 $n_i$ 为原子体密度，$\sigma_i$ 为该能级在光子能量 $h\nu$ 下的光电离截面，$\lambda_i$ 为电子在该基体中的非弹性平均自由程，$T(E_i)$ 为分析器透射函数，$D(E_i)$ 为探测器效率，$\theta$ 为相对表面法线的取出角，$E_i = h\nu - E_{\mathrm{B},i} - \varphi_{\mathrm{sp}}$ 为光电子动能。把与基体和仪器有关的因子归并成\emph{相对灵敏度因子}
\begin{equation}
S_i \;\propto\; \sigma_i(h\nu)\,\lambda_i(E_i)\,T(E_i)\,D(E_i)
\label{eq:ch09-rsf}
\end{equation}
则峰面积 $A_i \propto I_i$，原子分数为
\begin{equation}
C_i = \frac{A_i/S_i}{\sum_j A_j/S_j}
\label{eq:ch09-atomfrac}
\end{equation}
这就是灵敏度因子法。$\lambda$ 随动能缓慢增加（约 $E^{0.5\sim0.8}$；Seah--Dench 经验式给出以单层数表示的 $\lambda \approx 538E^{-2}+0.41a^{3/2}E^{1/2}$，$E$ 为动能／eV，$a$ 为单层厚度／nm），因此同一元素的不同谱线、不同元素之间的采样深度并不相同。
\emph{核对说明}：已对照原页面图像（书内 p.630、p.633 / PDF p.121、p.124）。原书给出式 (9-8)（含 $K\,T(\mathrm{KE})\,L_{ij}(\gamma)\,\sigma_{ij}$ 与深度积分）与半定量近似式 (9-10)（含 $\mathrm{EK}^{0.5}$），并说明单质材料在 100—2000 eV 间 $\lambda$ 近似正比于 $\mathrm{EK}^{0.5}$；本节的 $I_i \propto n_i\sigma_i\lambda_i T_i D_i\cos\theta$、$S_i$ 归并式与 $C_i$ 归一式为按学科标准形式的改写，与原书式 (9-8)、(9-9) 形式一致而符号与取舍不同。Seah--Dench 经验式及其常数（538、0.41）属原书未给出：原书 §9.2.5—§9.2.10 的 OCR 全文（\texttt{ch09-b.txt}）中未出现该式，Seah--Dench 仅见于本分片文末参考文献 [10]（书内 p.732 / PDF p.223），正文并未引用，故常数与「$E^{0.5\sim0.8}$」一并记为本改编本按通用文献补入。
\end{principle}

\begin{measure}
半定量的可靠性几乎完全由"条件一致性"决定：同一批样品、同一天、同一 X 射线源与功率、同一通能（透射函数 $T(E)$ 随通能变化）、同一步长与驻留时间、同一取出角。定量必须用高分辨窄扫描，宽扫描只用于元素识别。样品要压成平整致密的片——粉末直接撒在胶带上会同时引入粗糙阴影与胶带污染。记录样品是否含氢：XPS 测不到 H，报出的 at\% 是"除 H 外归一"的结果，与 XRF／ICP 的质量分数不可直接比较。
\end{measure}

\begin{dataflow}
窄扫谱 $\to$ 能量校准 $\to$ 背景扣除 $\to$ 分峰定面积 $\to$ 灵敏度因子归一（式 \equref{eq:ch09-atomfrac}）$\to$ 原子分数 $\to$ 与体相技术比较 $\to$ 结论。每一步都要能报出不确定度：面积拟合误差、灵敏度因子不确定度、背景模型带来的系统偏差。
\end{dataflow}

\paragraph{为什么只能算半定量：四类不可消除的偏差}
\begin{itemize}
\item \textbf{基体效应}。灵敏度因子通常在特定标样、特定仪器条件下测得，$\sigma$、$\lambda$、$T$ 都与基体密度、电子结构、带隙以及仪器透镜模式与通能有关；文献中的灵敏度因子不能跨仪器直接搬用，同一台仪器改变通能，$T(E)$ 曲线改变，定量结果也随之变化。
\item \textbf{逃逸深度差异}。Al K$\alpha$ 下 C 1s（动能约 1200 eV）的 $\lambda$ 约 2—3 nm，而 Ni 2p（动能约 630 eV）只有约 1.0—1.5 nm。表面存在成分梯度（偏析、钝化层、碳覆盖层、载体包覆）时，各元素的"原子比"实际来自不同深度；改变取出角该比值就会变，说明它不是一个体相量。
\item \textbf{样品不均匀}。粉末、粗糙表面、多孔载体、岛状分布的活性组分都会破坏式 \equref{eq:ch09-quant-int} 中的 $\cos\theta$ 假设（阴影、遮蔽、孔壁再吸收），表观浓度一般偏低；活性组分以颗粒存在时，颗粒内部信号被自身衰减，会进一步低估。
\item \textbf{污染层}。吸附碳氢层与水／氧化层使出射电子按 $\exp[-d/(\lambda\cos\theta)]$ 衰减；动能低的谱线衰减更重，因此表面金属组分（Ni 2p、Co 2p、Fe 2p）比污染碳（C 1s，动能高）被压得更低，表观金属含量系统性偏低。污染层厚度随抽真空时间与样品历史变化，同一个"原子比"在一天之内就可能漂移。
\end{itemize}
此外还有三个放大误差的来源：化学态改变会改变峰形与有效截面；卫星峰与多重分裂若归类不当会使面积失真；峰重叠（如 Ru 3d$_{5/2}$ 与 C 1s、Sb 3d$_{5/2}$ 与 O 1s）必须先分峰再报数。

\paragraph{检出限与误差怎么报}
常规条件下 XPS 的检出限约为 0.1—1 at\%（均匀样品、强谱线、充分累加计数）；轻元素与弱谱线更差。低于检出限的结果只能写成"未检出（$<$检出限）"。误差报告的通行做法：\emph{同批样品之间比较}取 $\pm5\%\sim\pm10\%$；\emph{绝对原子分数}（跨仪器、跨实验室）取 $\pm20\%\sim\pm30\%$。并且必须同时写明三件事：是否把污染碳计入总量、是否做过深度或取出角修正、检出限是多少。

\paragraph{读谱实例：一个表面 Ni 含量被污染层压低的算例}
\begin{obs}
同一条件下的窄扫面积（示意值，计数$\cdot$eV）：$A(\mathrm{C\ 1s})=2.0\times10^{4}$，$A(\mathrm{O\ 1s})=4.0\times10^{4}$，$A(\mathrm{Ni\ 2p_{3/2}})=1.2\times10^{3}$。取相对灵敏度因子 $S(\mathrm{C\ 1s})=0.25$、$S(\mathrm{O\ 1s})=0.66$、$S(\mathrm{Ni\ 2p_{3/2}})=4.2$（相对 F 1s 的量级值），按式 \equref{eq:ch09-atomfrac} 归一得 C $\approx$ 57 at\%，O $\approx$ 43 at\%，Ni $\approx$ 0.2 at\%。
\emph{核对说明}：已对照原页面图像（书内 p.632—p.633 / PDF p.123—p.124）。原书印有两套灵敏度因子：表 9-14 为 Scofield 计算的理论元素灵敏度因子（基于 C 1s $=1$，用时须增加分析深度项 $\lambda$，通常取 $\mathrm{KE}^{0.6}$），表 9-15 为由 PHI 谱仪公司提供的实验灵敏度因子（基于 F 1s $=1$，须按分析器类型——半球形或筒镜形——乘峰动能修正），并写明「一般使用元素灵敏度因子法准确度优于 15\%，在同一仪器上使用标准样品测量准确度优于 5\%」。本节算例的 $S(\mathrm{C\ 1s})=0.25$、$S(\mathrm{O\ 1s})=0.66$、$S(\mathrm{Ni\ 2p_{3/2}})=4.2$ 与表 9-15 同以 F 1s 为基准，量级一致，但未逐值照抄原书表（页面图像分辨率不足以保证逐值读取），属按通用文献量级取的示意值；按式 \equref{eq:ch09-atomfrac} 复算 $A/S$ 依次为 $8.0\times10^{4}$、$6.06\times10^{4}$、$2.86\times10^{2}$，归一后 C $\approx57$ at\%、O $\approx43$ at\%、Ni $\approx0.20$ at\%，与本例所报数值自洽。实际报数须改用本台仪器、本通能下测得的灵敏度因子。
\end{obs}
\begin{infer}
不能由 0.2 at\% 直接推断"Ni 流失或分散度极低"：若表面存在约 2 nm 的碳质污染层，Ni 2p（动能约 630 eV，$\lambda\approx1.3$ nm）的衰减因子约 0.2，而 C 1s 与 O 1s（$\lambda\approx2.5$—3 nm）约 0.5，表观 Ni 分数会被压低约 2—3 倍。另外 0.2 at\% 已接近检出限，相对不确定度很大。\emph{处置}：用 \ce{Ar+} 轻刻蚀（注意择优溅射与离子束还原）或在较大取出角下复测，报数时注明"含污染层、未做深度修正"。
\end{infer}

\begin{pitfall}
（1）把 at\% 当准确值，不报误差、不报检出限（见 §\ref{sec:ch09-problems}）。（2）直接引用文献灵敏度因子而不核对仪器与通能。（3）把污染碳计入组成，据此讨论"积炭"或"碳覆盖度"——必须先证明那部分碳属于样品（峰形、结合能、$\pi$-$\pi^{*}$ 伴峰、刻蚀前后对比）。（4）用动能相差很大的两根谱线比较"表面"与"体相"而不考虑 $\lambda$ 差异。（5）把 \ce{Ar+} 刻蚀后的定量结果当作原始表面：择优溅射与离子束诱导的还原、偏析都会改变表层组成。（6）把灵敏度因子归一后的表面浓度直接与催化活性关联，而不先证明该表面物种就是活性位。
\end{pitfall}

\begin{crosscheck}
体相组成用 XRF／ICP-AES 回查；最表层（单原子层）组成用 LEIS 与 RBS 回查；微区组成用 EDS／WDS 回查（采样深度在微米量级，与 XPS 的纳米量级恰好互补）。在同一台仪器上用同基体标样做工作曲线，可以把半定量提升为"带标样的定量"。
\end{crosscheck}

\subsection{XPS 数据的后期处理}
\label{sec:ch09-postproc}

\emph{核对说明}：本节有原书对应正文，已对照原页面图像（书内 p.633—p.635 / PDF p.124—p.126）并通查其 OCR 全文（\texttt{sources/parts/ch09-b.txt} 第 589—779 行，即原书 §9.2.9「XPS 数据的后期处理」，书内 p.633 起）。原书按「1. 荷电校正／2. 使用非单色化激发源，需进行卫星峰扣除／3. 光电子峰的识别，定性分析／4. 适度的谱线平滑处理／5. 本底扣除（线性与非线性两种；shirley 本底、smart 本底、tougaard 本底）」的顺序列出，并对本底成因（非弹性散射电子与次级电子在光电子谱峰的高结合能侧形成本底，轫致辐射本底随光电子动能降低而增强）与 shirley 本底的定义作了文字说明，但全节未给出任何背景扣除公式。本节正文（含式 \equref{eq:ch09-shirley}、\equref{eq:ch09-tougaard} 与分峰约束）为按学科标准形式的改写：方法分类与原书一致，公式、符号与分峰约束属改编补充（原书未印）。

\begin{probout}
后处理只回答三个问题：能量标尺可信吗（荷电校准）；峰下面的背景怎么去掉（面积才可比）；峰能拆成几个化学态、这个拆法唯一吗（背景模型、线型约束与不确定度）。这三步中任何一步没有交代清楚，后面报出的"化学态比例"都无法复核。
\end{probout}

\begin{principle}
\textbf{背景的物理来源}\quad 光电子在逸出途中发生非弹性散射，损失能量后落到主峰\emph{高结合能侧}，形成阶跃状背景；叠加的二次电子级联则给出随结合能缓慢上升的宽背景。所以 XPS 谱的背景总是"往高结合能方向抬升"。

\textbf{线性背景}\quad 取峰两侧平坦区的两点连成直线，假设峰区背景近似线性。算法最简单、结果可复现，适合窄能区、背景平坦、两端有明确平坦区的场合；它的风险是把台阶状背景当成直线，在台阶明显时系统性过扣高结合能一侧。

\textbf{Shirley 背景}\quad 假设任一点处的非弹性背景正比于该点\emph{低结合能侧}已经积累的峰强度（即所有可能在此处贡献损失电子的源），迭代求解：
\begin{equation}
B^{(n+1)}(x) = k\int_{x_0}^{x}\bigl[\,I(x')-B^{(n)}(x')\,\bigr]\,\mathrm{d}x'
\label{eq:ch09-shirley}
\end{equation}
其中 $k$ 由两端点的背景取值确定。它对台阶状背景的适应性最好，是目前最常用的做法；代价是结果依赖端点选择、需要迭代收敛（有时不收敛），而且它本身不带物理参数，不能用来反推损失函数。

\textbf{Tougaard 背景}\quad 用电子非弹性散射的普适截面 $K(T)$（$T$ 为能量损失）卷积全谱得到背景：
\begin{equation}
B(E) = \int_{E}^{E_{\max}} \bigl[\,I(E')-B(E')\,\bigr]\,K(E'-E)\,\mathrm{d}E'
\label{eq:ch09-tougaard}
\end{equation}
普适参数常取 $K(T)\propto BT/(C+T^{2})^{2}$，$B\approx2866\ \mathrm{eV^{2}}$、$C\approx1643\ \mathrm{eV^{2}}$。它是有物理依据的背景模型，适合金属（Cu、Au、Ag 等）中显著的等离子体损失与宽损失分布，并可把损失特征一并纳入背景；对绝缘体、复杂氧化物与窄能区的孤立单峰，参数不再具有普适性，容易过扣。
\emph{核对说明}：已对照原页面图像（书内 p.634—p.635 / PDF p.125—p.126）。原书 shirley 本底写为「这样给定的动能处的本底，是与本底以上较高动能处的谱峰总面积成固定比例」，即任一点背景正比于其高动能（低结合能）侧已积累的峰强度——式 \equref{eq:ch09-shirley} 的积分方向与此一致（原书只作文字说明、未印积分式），判定为形式一致。tougaard 本底原书亦只给文字定义（「用关于相对于能量损失的微分非弹性散射截面模型和表面区发射原子的三维分布模型所得到的强度分布」，\texttt{sources/parts/ch09-b.txt} 第 671 行）；普适参数 $B\approx2866\ \mathrm{eV^{2}}$、$C\approx1643\ \mathrm{eV^{2}}$ 原书未给出（该文件全文未出现这两个常数），系按通用文献补入，属原书未给出而非原书印错。
\end{principle}

\begin{table}[tbp]
\centering
\caption{三种背景扣除方法的假设、适用范围与主要风险。选背景的依据是背景的物理形态（台阶是否明显、有无显著等离子体损失），不是拟合优度。\src{本表为本改编本归纳，非源书表格；本节源文件中无对应 OCR 文本，三种背景方法的适用范围与风险按学科标准形式归纳，待回查原书。}}
\label{tab:ch09-bg}
{\fontsize{8}{9.6}\selectfont
\begin{tabular}{@{}p{15mm}p{44mm}p{34mm}p{46mm}@{}}
\toprule
方法 & 假设与算法 & 适用 & 主要风险 \\
\midrule
线性 & 峰两侧平坦区两点连线，假设峰区背景为直线；无参数 & 窄能区、背景平坦、两端有明确平坦区 & 台阶背景下系统性过扣高结合能一侧；端点选在峰尾时面积偏低一到三成 \\
Shirley & 背景正比于低结合能侧累积的峰强度，迭代求 $k$（式 \equref{eq:ch09-shirley}） & 最通用；台阶状背景的金属、氧化物、宽能区扫描 & 依赖端点选择；迭代可能不收敛；无物理参数，不能反推损失函数；容易把弱真实信号当背景吃掉 \\
Tougaard & 用非弹性散射普适截面卷积全谱（式 \equref{eq:ch09-tougaard}） & 金属中显著的等离子体损失与宽损失分布；可同时给出损失结构 & 参数对绝缘体与复杂氧化物不普适；窄能区单峰容易过扣；需额外参数与判据 \\
\bottomrule
\end{tabular}}
\end{table}

\begin{measure}
（1）步长应不大于最小半高宽的 1/10，否则面积的内插误差会进入定量。（2）背景端点必须落在真正的平坦区，最好在峰两侧等距离处；端点稍有移动，面积可变化百分之十几。（3）荷电校准所用的参考峰与中和枪设置必须与采集同时记录，参考峰与样品峰最好来自同一次装样、同一条件。（4）拟合时必须记录：所用线型、半高宽约束、自旋轨道分裂与面积比约束、以及每个峰的拟合标准误差。
\end{measure}

\begin{dataflow}
原始谱 $\to$（1）能量标定（荷电校准）$\to$（2）背景模型与端点选择 $\to$（3）线型与约束设置 $\to$（4）拟合 $\to$（5）面积比 $\to$（6）化学态比例及其不确定度 $\to$（7）与独立证据对照。误差按此顺序传递：荷电校准决定峰位归属，背景决定面积，约束决定组分比——后一步无法修正前一步的系统错误。
\end{dataflow}

\paragraph{荷电校准：污染碳 284.8 eV 与 Au 4f$_{7/2}$}
\begin{itemize}
\item \emph{污染碳法}：把 C 1s 中的污染碳组分定在 284.8 eV，其余峰位整体平移。前提是样品本体不含碳、污染层与样品荷电相同、且污染碳确实是 sp$^{3}$ 型的 C—H／C—C。
\item \emph{Au 4f$_{7/2}$ 法}：在样品上蒸镀或贴附金属 Au（或让 Au 箔与样品电接触），把 Au 4f$_{7/2}$ 定在 84.0 eV。前提是 Au 与样品电接触良好（否则标定的是 Au 自己的电位），并且 Au 颗粒足够大以避免尺寸效应、不与样品发生相互作用。
\item \emph{首选做法}：用样品自身化学态已知的物种做内标（如氧化物晶格氧 O 1s 529.8—530.2 eV、金属态的费米边），同时用两个独立内标比对。若两者给出不同的平移量 $\Delta$，说明存在差分荷电（见 §\ref{sec:ch09-problems}），此时不能只取其中一个。
\item \emph{陷阱}：污染碳的结合能并非恒定（文献报道在 284.3—285.2 eV 之间），污染层本身还含 C—O 与 C=O 组分；样品本体含碳时（石墨 284.6 eV、\ce{C60} 284.75 eV、碳化物 282—283 eV、聚合物 285.0 eV、\ce{C-O} 286.5 eV、碳酸盐 290.3 eV），污染碳峰与样品碳峰重叠，"把最强峰调到 284.8 eV"会连样品碳峰一起搬走。校正的方向与量级必须在正文写明，并给出所用参考与不确定度（通常 $\pm0.2$ eV）。
\end{itemize}

\paragraph{分峰拟合的约束}
线型：绝缘体与聚合物用 Gaussian--Lorentzian 混合 $\mathrm{GL}(m)$（$m$ 为 Lorentz 分数，常用 GL(30)）；金属因电子-空穴对激发带有向高结合能侧的拖尾，需用 Doniach--Šunjić 线型或带不对称参数的 GL。约束清单如下。
\begin{itemize}
\item \emph{峰位}：由化学态决定，必须与同批样品的其他谱线、与文献或标样一致，不允许为改善拟合而任意移动。
\item \emph{半高宽}：同一元素各化学态的半高宽差异一般不超过 0.2—0.3 eV（\ce{C\ 1s} 污染碳在单色 Al K$\alpha$ 下通常为 1.2—1.6 eV）。若必须让某个组分有异常大的半高宽才能收敛，说明模型可疑（可能把差分荷电或未解析的多个物种压进了一个峰）。
\item \emph{自旋轨道双峰}：分裂 $\Delta$ 由元素决定，拟合时应固定或只允许极小浮动；面积比由简并度 $(2j+1)$ 决定：p 为 2:1（$2p_{3/2}:2p_{1/2}$），d 为 3:2（$3d_{5/2}:3d_{3/2}$），f 为 4:3（$4f_{7/2}:4f_{5/2}$）。常用值见表 \tabref{tab:ch09-so}。
\item \emph{峰数}：不超过数据能支持的数目（每个组分一般要有可辨认的肩峰、独立的谱线约束或化学合理性支持）；拟合完成后必须报 $\chi^{2}$ 或残差与每个峰面积的误差。
\end{itemize}

\begin{table}[tbp]
\centering
\caption{常用芯能级的自旋轨道分裂与强度比，以及拟合时的注意事项。强度比由简并度 $(2j+1)$ 决定，与实验值的偏差本身就是一个诊断量。\src{本表为本改编本归纳，非源书表格；「S 2p」一行据 S2 表 9-37 的实测统计，书内 p.712 / PDF p.203，其余分裂值与化学位移范围按通用文献量级给出。}}
\label{tab:ch09-so}
{\fontsize{8}{9.6}\selectfont
\begin{tabular}{@{}p{16mm}p{24mm}p{18mm}p{56mm}@{}}
\toprule
谱线 & 分裂 $\Delta$/eV & 强度比 & 拟合时的约束与注意 \\
\midrule
S 2p & 1.16—1.26（平均 1.19） & 2:1（实测平均 1.91） & 取自源文件表 9-37 的实测统计，可作约束的可信起点 \\
P 2p & 0.84—0.87 & 2:1 & 常与相邻峰重叠，须固定 $\Delta$ 与强度比 \\
Cl 2p & 1.6 & 2:1 & 有机氯污染的判据；与 Si 2s 区接近时注意重叠 \\
Ti 2p & 5.7 & 2:1 & 多个价态的 $2p_{3/2}$ 相隔 1.5—1.8 eV，易与展宽或差分荷电混淆，需卫星峰辅助 \\
Cu 2p & 19.8 & 2:1 & \ce{Cu^{2+}} 有强卫星峰（约 940—945 eV），\ce{Cu^0} 与 \ce{Cu+} 无；卫星峰不能计入主峰面积 \\
Ni 2p & 17.3 & 2:1 & \ce{Ni^0} 与 \ce{Ni^{2+}} 的 $2p_{3/2}$ 相差约 1.5—2 eV，卫星结构不同；金属 Ni 需不对称线型 \\
Mo 3d & 3.15 & 3:2 & \ce{Mo^{6+}} 与 \ce{Mo^{4+}} 相差约 1.6—1.8 eV；S 2s（约 226 eV）与 Mo 3d 重叠，必须一并拟合 \\
Pd 3d & 5.26 & 3:2 & \ce{Pd^0} 与 \ce{Pd^{2+}} 相差约 1.5—2 eV；高分散 Pd 峰形不对称，慎用对称线型 \\
Pt 4f & 3.33 & 4:3 & \ce{Pt^0}、\ce{Pt^{2+}}、\ce{Pt^{4+}} 依次向高结合能移动约 1—2 eV \\
Au 4f & 3.67 & 4:3 & 常作荷电校准内标（$4f_{7/2}=84.0$ eV）；同时核对 $4f_{5/2}$ 位置与峰形 \\
\bottomrule
\end{tabular}}
\end{table}
\emph{核对说明}：已对照原页面图像（书内 p.712 / PDF p.203）核对「S 2p」一行：源文件表 9-37（硫元素在聚合物中 S 2p 结合能值）列出 PETH、PHMS、PPS、PES、PSS 五个聚合物，$\Delta$ 依次为 1.19、1.26、1.18、1.16、1.17 eV（平均 1.19），面积比 $I(2p_{3/2})/I(2p_{1/2})$ 依次为 1.84、1.91、1.90、1.84、2.04（平均 1.91），与表中数值一致，判定为已核对一致。其余各行（P 2p、Cl 2p、Ti 2p、Cu 2p、Ni 2p、Mo 3d、Pd 3d、Pt 4f、Au 4f）的分裂值与化学位移范围按通用文献量级给出：原书（书内 p.694—733）的表格均为 C 1s、O 1s、N 1s、S 2p 等结合能统计（表 9-30—表 9-37），未给上述其他芯能级的自旋轨道分裂值，属原书未给出而非印错。

\begin{figure}[tbp]
\centering
\begin{tikzpicture}
\tikzset{declare function={
  lor(\x,\a,\w)=1/(1+((\x-\a)/\w)*((\x-\a)/\w));
  pA(\x)=30*lor(\x,5,0.8);
  pB(\x)=18*lor(\x,7,0.9);
  cum(\x)=(24*(90-atan((\x-5)/0.8))+16.2*(90-atan((\x-7)/0.9)))*0.0174533;
  bsh(\x)=5+0.1188*cum(\x);
  blin(\x)=20.34-0.8933*\x;
  spec(\x)=pA(\x)+pB(\x)+bsh(\x)+0.25*\x;
}}
\begin{groupplot}[
  group style={group size=1 by 3, vertical sep=8mm},
  width=\onecolwidth,
  height=31mm,
  xmin=0, xmax=12, ymin=0,
  samples=241,
  tick label style={font={\fontsize{8}{9.6}\selectfont}},
  label style={font={\fontsize{8}{9.6}\selectfont}},
  axis line style={line width=0.5pt},
  every axis plot/.append style={line width=0.6pt},
]
\nextgroupplot[ymax=56, ylabel={计数（示意）}, ytick={0,20,40}]
\addplot[black, domain=0:12] {spec(x)};
\addplot[black, dashed, domain=0:12] {blin(x)};
\addplot[black, dotted, domain=0:12] {bsh(x)};
\node[align=left, text width=44mm, anchor=north west,
      font={\fontsize{8}{9.6}\selectfont}] at (rel axis cs:0.03,0.97)
      {实线：模型谱\\ 虚线：线性背景\\ 点线：Shirley 背景};
\nextgroupplot[ymax=40, ylabel={计数（示意）}, ytick={0,15,30}]
\addplot[black, domain=0:12] {spec(x)-blin(x)};
\addplot[black, dashed, domain=0:12] {pA(x)+pB(x)};
\addplot[gray, dotted, domain=0:12] {pA(x)};
\addplot[gray, dotted, domain=0:12] {pB(x)};
\node[align=left, text width=52mm, anchor=north east,
      font={\fontsize{8}{9.6}\selectfont}] at (rel axis cs:0.97,0.97)
      {线性背景扣除后\\ 残差 $+1.8$（4 eV）$\to$ $-4.0$（8 eV）\\ 符号反转：必须改组分比};
\nextgroupplot[ymax=40, xlabel={结合能（相对标度）}, ylabel={计数（示意）}, ytick={0,15,30}]
\addplot[black, domain=0:12] {spec(x)-bsh(x)};
\addplot[black, dashed, domain=0:12] {pA(x)+pB(x)};
\addplot[gray, dotted, domain=0:12] {pA(x)};
\addplot[gray, dotted, domain=0:12] {pB(x)};
\node[align=left, text width=52mm, anchor=north east,
      font={\fontsize{8}{9.6}\selectfont}] at (rel axis cs:0.97,0.97)
      {Shirley 背景扣除后\\ 残差为 $0\to3$ 的平缓斜坡\\ 不造成组分比的系统偏差};
\end{groupplot}
\end{tikzpicture}
\caption{模型计算：同一张谱分别用线性背景（中）与 Shirley 背景（下）扣除后的双峰分解差异。读者应比较两个残差：中图残差在 4 eV 附近为 $+1.8$、在 8 eV 附近为 $-4.0$，\emph{出现符号反转}，说明背景拿走了本属于峰或留下了本属于背景的强度，且高低结合能端方向相反，无法靠微调峰面积吸收，只能改变两个组分的比例才能收敛；下图残差只是 $0\to3$ 的单调斜坡，不会造成组分比的系统偏差。图内实线为谱，粗虚线为固定的双峰模型和，细点线为两个组分。条件与局限：两个 Lorentz 型组分（峰位 5 与 7、半高宽 1.6 与 1.8、峰高 30 与 18）叠加台阶背景（阶高约 15）与缓升二次电子背景，未含真实谱的卫星峰、多重分裂与统计噪声；在本示意参数下改用线性背景使高结合能组分的表观相对含量下降约 25\%—30\%。\modelmark\figsrc{本图为模型计算，源文件无对应图}}
\label{fig:ch09-bgfit}
\end{figure}

\paragraph{读图实例：背景模型如何改变组分比}
\begin{obs}
图 \figref{fig:ch09-bgfit} 中，模型谱由两个 Lorentz 型组分叠加台阶背景与缓升的二次电子背景构成。用线性背景扣除后，与"固定峰位、固定半高宽的双峰模型"相减得到的残差在低结合能端为正（约 $+1.8$）、在高结合能端为负（约 $-4.0$），出现符号反转；用 Shirley 背景扣除后，残差只是一条 $0\to3$ 的平缓单调斜坡，没有符号反转。
\end{obs}
\begin{infer}
残差符号反转意味着背景模型在高、低结合能端分别多扣与少扣了强度，方向相反，无法靠微调某一峰的面积同时吸收，只能\emph{改变两个组分的比例}才能收敛。也就是说：同一张谱、同一组峰位与半高宽，仅更换背景模型，就会得到不同的化学态比例。在本示意参数下，改用线性背景使高结合能组分的表观相对含量下降约 25\%—30\%。
\emph{核对说明}：已对照原页面图像（书内 p.634—p.635 / PDF p.125—p.126）。原书讨论本底扣除时只作定性比较：线性扣除法适合纯金属样品、方法简单可操作性强但缺乏理论依据，非线性扣除法适合所有样品（书内 p.634）；图 9-79 以 Fe 2p 为例给出选点不同造成的峰面积差——线性本底 1 点峰面积 44560、2 点 66589（差 22029），非线性本底 1 点 41167、2 点 49814（差 8647），原书并未把它折算为组分比偏差，也没有本节这样的双峰模型算例。故本处「15\%—20\%」由本图模型参数（两 Lorentz 组分的峰位 5 与 7、半高宽 1.6 与 1.8、峰高 30 与 18，台阶背景阶高约 15）解析估算得出，未做数值积分，属本改编本的模型估算，原书无同类算例可对照。
\end{infer}

\begin{pitfall}
（1）把分峰结果当唯一解：换背景、换端点、换 GL 比例、换峰数都能得到"看起来一样好"的拟合（$\chi^{2}$ 差不到 1\%，组分比却可能差 30\%）。（2）背景端点选在峰的尾部而不是平坦区，系统性过扣。（3）固定了半高宽却不固定自旋轨道分裂与面积比，造出并不存在的物种。（4）用 C 1s 284.8 eV 标定本身含碳的样品（见 §\ref{sec:ch09-problems}）。（5）用未扣背景的原始谱直接比较峰高——峰高受半高宽与背景影响，只有面积才有物理意义。（6）只看拟合优度，不看化学合理性。
\end{pitfall}

\begin{crosscheck}
（1）用同一元素的另一根谱线（Ni 2p 与 Ni 3p、Cu 2p 与 Cu 3p）独立拟合，两组分比应一致；（2）用 Auger 参数交叉验证化学态归属；（3）换两种背景模型、两种端点重复拟合，组分比变化应小于 10\%，否则只能报区间；（4）用不同取出角或不同动能谱线检查是否存在深度不均；（5）与 XANES／EXAFS、Raman、\ce{H2}-TPR 等独立技术对照。
\end{crosscheck}

\subsection{XPS 实验中的常见问题}
\label{sec:ch09-problems}

\emph{核对说明}：本节有原书对应正文，已对照原页面图像（书内 p.640—p.641 / PDF p.131—p.132）并通查其 OCR 全文（\texttt{sources/parts/ch09-b.txt} 第 780 行起，即原书 §9.2.10「XPS 分析中常遇到的一些问题」，书内 p.640 起）。原书列举的问题为：XPS 的检出范围、信号强度与分辨能力、谱仪的能量标尺校准、XPS 的伴峰、绝缘样品表面的荷电校正、各种谱线的重叠、对被测样品的要求、实验中 X 射线对样品表面的副效应等，并按「1. XPS 的检出范围」「2. 信号强度与分辨能力」逐条展开。原书用「检出范围」而不用「检出限」，并在脚注中写明「有种提法为所谓的『检出限』，这种提法不妥。实际上无检出极限，因为辐射源的功率可再提高，检测器的灵敏度也可再提高」；早期仪器相对灵敏度约 1\%（相对原子百分比），现在可达约千分之一。本节把常见问题重组为荷电、污染、X 射线损伤、差分荷电、样品制备假象五类并按「症状—判据—处置」列表，属按学科标准形式与仪器惯例的改编补充，与原书的分类与叙述顺序不同；数量级（约 0.1—1 at\%）与原书一致。

\begin{probout}
把"样品问题"与"样品本征性质"分开，是 XPS 实验设计的第一要务。以下五类问题覆盖了绝大多数误判：荷电、污染、X 射线损伤、差分荷电、样品制备引入的假象。每一类都按"症状 $\to$ 判据 $\to$ 处置"给出可操作流程（表 \tabref{tab:ch09-problems}）。
\end{probout}

\begin{table}[tbp]
\centering
\caption{XPS 实验中的常见问题：症状、判据与处置。判据一栏给的是可以在谱图上直接核对的量，凡判据不成立就不要把现象写成结论。\src{本表为本改编本归纳，非源书表格；本节源文件中无对应 OCR 文本，五类问题的症状、判据与处置按学科标准形式与仪器惯例归纳，待回查原书。}}
\label{tab:ch09-problems}
{\fontsize{8}{9.6}\selectfont
\begin{tabular}{@{}p{15mm}p{36mm}p{40mm}p{48mm}@{}}
\toprule
问题 & 典型症状 & 判据（可直接核对） & 处置 \\
\midrule
荷电（整体） & 全部峰向高结合能整体平移 1—10 eV，峰形不变；平移量随 X 射线通量与时间变化 & 降低通量或开启中和枪后\emph{所有}峰同步回移；各元素峰的位移量一致 & 降低通量、开中和枪、改善与样品台的接触（压薄、压实在台上、金箔搭接）；用内标校准并写明所用参考与 $\pm0.2$ eV 的不确定度 \\
表面污染 & C 1s 与 O 1s 随时间增强；表面金属组分（Ni 2p、Co 2p）相对含量下降；C 1s 在 286—289 eV 出现组分 & 抽真空时间越长污染越重；新解理或刚刻蚀后立即测量则显著减弱；空白样品台与胶带在同条件下出现同样的 C／O & 缩短样品在分析室的停留时间；先低分辨快速扫一遍再定条件；必要时加热除气或轻刻蚀（会改变表面状态，须对照）；报数时说明污染层是否计入总量 \\
X 射线损伤 & 低价态随时间单调增长（\ce{Cu^{2+}}$\to$\ce{Cu+}、\ce{Mo^{6+}}$\to$\ce{Mo^{4+}}、\ce{Ce^{4+}}$\to$\ce{Ce^{3+}}）；聚合物脱氟、C—O 组分增强、O 1s 增强 & 同一位置连续采集 8—10 次，价态比或峰形随累积剂量\emph{单调}变化；移动到未照射的新位置后恢复 & 降低 X 射线功率、用单色器、缩短单次采集多点平均、低温台；先测最敏感的元素；用剂量序列外推到零剂量 \\
差分荷电 & 同一元素的峰展宽或分裂，间距随通量／时间变化；样品峰与污染碳峰的漂移量不一致 & 分裂间距与已知化学位移不符；降低通量或改善电接触后分裂消失；样品峰与内标峰给出的 $\Delta$ 不同 & 重新制样（压薄压实、导电胶或金箔搭接）、降低通量、适当中和；用第二根谱线（如 Ni 3p）与双内标验证后重测 \\
制备假象 & 出现与制备过程有关的元素（Si、Na、F、S、增塑剂）；刻蚀后出现 Ar、低价态或组分突变；颗粒表面出现异常氧化物 & 空白对照（胶带、溶剂、模具、\ce{Ar+} 刻蚀）在同条件下测出相同指纹；未处理的同批样品无此信号 & 改用无胶带固定（压片、金属夹）、换溶剂并充分干燥、避免过度研磨；刻蚀条件先用已知厚度标样标定速率并做空白；完整记录样品历史 \\
\bottomrule
\end{tabular}}
\end{table}

\begin{principle}
\emph{荷电}：光电离不断把电子送出样品，绝缘样品表面累积正电荷，稳态电位由 X 射线通量、样品电导、与样品台的接触、二次电子与中和电子的供给共同决定，因此它随时间、位置与通量漂移（表现为整体平移，不改变峰形）。\emph{污染}：真空中残余的 \ce{H2O}、CO 与碳氢化合物吸附在表面，并被 X 射线与二次电子诱导分解，形成 1—3 nm 的碳质层，厚度随抽真空时间增长。\emph{X 射线损伤}：二次电子与 Auger 过程造成键断裂与脱附——聚合物交联或脱氟，金属氧化物被还原（\ce{Cu^{2+}}$\to$\ce{Cu+}、\ce{Mo^{6+}}$\to$\ce{Mo^{4+}}、\ce{Ce^{4+}}$\to$\ce{Ce^{3+}}），表现为"低价态"随时间增长。\emph{差分荷电}：同一视场内不同区域、不同颗粒、污染层与基体之间的电导不同，各自建立不同的稳态电位，于是同一元素的谱线出现\emph{随通量与时间变化的}展宽或分裂——它不是化学位移。\emph{制备假象}：研磨引入氧化物与相变，胶带与脱模剂引入 Si、Na、增塑剂，溶剂残留引入 C—O、F、S，\ce{Ar+} 刻蚀引入择优溅射、离子束还原、Ar 植入与表面粗糙化。
\end{principle}

\begin{measure}
诊断性实验设计（按代价从低到高）：
\begin{enumerate}
\item \emph{时间序列}：固定位置连续采集 8—10 次同一窄扫，作峰位、半高宽、组分比随累积剂量的变化曲线。峰位单调漂移或低价态单调增长 $\Rightarrow$ 损伤或荷电；峰形不变而只整体平移 $\Rightarrow$ 荷电。
\item \emph{通量序列}：把 X 射线功率降低 3—5 倍重测。峰位随通量变化 $\Rightarrow$ 荷电；峰位不变而低价态比例明显下降 $\Rightarrow$ 损伤。
\item \emph{中和枪开／关对照}：只影响荷电，不影响污染与损伤。
\item \emph{位置对照}：移到未照射过的新位置立即采集。若"低价态"消失 $\Rightarrow$ 是束流效应而不是样品本征。
\item \emph{取出角序列}（15°／45°／90°）：判断污染层与深度分布，也用于验证"表面富集"是否真在表面。
\item \emph{空白对照}：样品台、胶带、溶剂、压片模具在同条件下各测一次，建立污染指纹库。
\item \emph{记录样品历史}：暴露空气时长、研磨与压片方式、溶剂与粘接方式，以及本次测量的时刻。
\end{enumerate}
\end{measure}

\begin{dataflow}
症状 $\to$ 判据（可量化的阈值，见表 \tabref{tab:ch09-problems}）$\to$ 处置 $\to$ 复测确认（问题是否消失或减轻）$\to$ 写结论时必须同时写明限制条件（是否含污染层、是否可能有损伤、是否做过荷电校准及所用参考）。
\end{dataflow}

\paragraph{读谱实例：看着像两个价态，其实是差分荷电}
\begin{obs}
Ni 2p$_{3/2}$ 区主峰半高宽由正常的约 1.4 eV 展宽到 3.2 eV，峰形不对称；强行拟合成两个峰得到 855.2 eV 与 857.4 eV，间距 2.2 eV，面积比约 1.5:1。连续采集 6 次后主峰位置由 855.2 eV 漂到 856.1 eV，而同时记录的污染碳 C 1s 只漂移 0.1 eV；把 X 射线功率降低 5 倍后，峰位回移约 0.6 eV，峰形变窄。
\end{obs}
\begin{infer}
三条判据都指向\emph{差分荷电}，而不是"两个价态"：（a）Ni 2p 的化学位移通常在 1—1.5 eV 量级，且不同价态伴随不同的卫星峰结构，2.2 eV 的分裂与常见价态对不上；（b）化学态比例不应随 X 射线功率变化；（c）样品峰与污染碳峰的漂移量不一致，说明两者处在不同电位上。\emph{处置}：重新制样（压薄、压实在样品台上，边缘用导电胶或金箔搭接）、降低 X 射线通量、适当中和后重测；并用同一元素的第二根谱线（Ni 3p）与双内标（C 1s 与 Au 4f$_{7/2}$）比对平移量来验证。
\end{infer}

\begin{pitfall}
四条最常见的误判，以及各自的判据与处置：
\begin{enumerate}
\item \textbf{把分峰结果当唯一解}。判据：换背景、换端点或换 GL 比例后组分比变化超过 10\%。处置：做敏感性分析，报组分比的区间而不是单值，并在正文写明约束条件。
\item \textbf{用 C 1s 284.8 eV 标定，而样品本身就有碳物种}。判据：样品合成或组成中含碳（碳化物、石墨、\ce{C60}、聚合物、有机配体、碳酸盐），或 C 1s 峰形不对称、半高宽大于 2 eV、在 286—290 eV 还有组分。处置：改用样品自身已知物种（晶格氧、金属态、费米边）或 Au 4f$_{7/2}$ 内标，或者明确按"未做荷电校准、只报相对位移"处理，并把所用参考写进正文。
\item \textbf{把差分荷电当成多价态}。判据：分裂间距与已知化学位移不符、峰位随通量／时间／位置变化、样品峰与内标峰漂移不一致、改善电接触或降低通量后分裂消失（见上面的读谱实例）。
\item \textbf{定量不报误差与检出限}。判据：只给"某元素 $=$ x at\%"而不给不确定度、检出限与是否计入污染碳。处置：报原子分数 $\pm$ 相对不确定度（同批比较 $\pm5\%\sim\pm10\%$，绝对 $\pm20\%\sim\pm30\%$），注明检出限（约 0.1—1 at\%）与归一化范围（不含 H）。
\end{enumerate}
此外还有两条同样常见：把 shake-up 伴峰或 X 射线卫星线当独立化学态定量；把 \ce{Ar+} 刻蚀后的结果当作原始表面状态。
\end{pitfall}

\begin{crosscheck}
XRD 回查体相物相；XANES／EXAFS 回查价态与配位（且可在原位气氛下测量，避免样品离开气氛后重构）；Raman 与 IR 回查碳物种与官能团；\ce{H2}-TPR 回查可还原性。最有力的一个交叉检验是\emph{换激发源}：芯能级结合能不随光子能量变化，而 Auger 峰动能不变（于是它的表观结合能随 $h\nu$ 平移）。若一个"化学态"的位置随激发源改变，它多半是 Auger 峰或卫星线的误判。
\end{crosscheck}

\srcnote{本分片按任务书覆盖 §9.2.7—§9.2.10。价带谱与 shake-up 伴峰部分取自源文件 \texttt{ch09-d.txt} 的 §9.6.2（书内 p.705—708 / PDF p.196—199，已对照页面图像）及表 9-37（书内 p.712 / PDF p.203，已对照页面图像）；§9.2.7—§9.2.10 的原书正文本身见书内 p.623—641，其 OCR 文本在另一分片的源文件 \texttt{sources/parts/ch09-b.txt}（§9.2.7 第 227 行／§9.2.8 第 415 行／§9.2.9 第 589 行／§9.2.10 第 780 行）。定量、后处理与常见问题三节的公式、常数与分类为按学科标准形式与仪器惯例的改写：与原书形式一致之处、原书未给出或未印之处，逐条记于各节前的核对说明。}
